\section{Выбор \latGCKgen\ с двенадцатью соседями}
\label{sec:fcc-choice}

\begin{definition}[Типы упаковки в трёх измерениях]
\label{def:packing-types}
Сравнивая носители, используем стандартную кристаллографическую терминологию
(без англоязычных аббревиатур SC, BCC, FCC, HCP в прозе):
\begin{center}
\small
\begin{tabular}{@{}p{0.52\linewidth}cc@{}}
\toprule
название & сокр. & $|N|$ \\
\midrule
\latPK & \latPKshort & $6$ \\
\latOCK & \latOCKshort & $8$ \\
\latGCK & \latGCKshort & $12$ \\
\latGPU & \latGPUshort & $12$ \\
\bottomrule
\end{tabular}
\end{center}
Носитель модели --- \latGCK\ $\carrierLattice\subset\mathbb{R}^3$.
В формулах иногда пишут $\mathrm{FCC}$ как условное обозначение этого типа Браве.
\end{definition}

\begin{theorem}[Носитель в $(3{+}1)$]
\label{thm:fcc-choice}
Пусть выполнены: замощение без дыр; единственная длина~$\caStepSpace$;
только светоподобные рёбра за один~$\htick$;
изотропия вакуума; один носитель на узел; физическое пространство трёхмерно.
Тогда $\carrierLattice$ --- \latGCK, $|N(x)|=12$ для всех узлов,
$\kappa_{\mathrm{link}}=1/12$.
\end{theorem}

\begin{proof}
Сначала отбираем упаковки с одной длиной ребра и без дыр, затем отсекаем вторую оболочку световым конусом.

\emph{Упаковка.}
\pusing{модели шара радиуса $\caStepSpace/2$ на узле}
Ячейку представляем шаром радиуса $\caStepSpace/2$: касание двух шаров есть ребро длины~$\caStepSpace$.
В трёх измерениях плотнейшая упаковка равных шаров --- гранецентрированная кубическая или гексагональная плотная (теорема Кеплера).
У каждого шара тогда ровно двенадцать соседей на расстоянии~$\caStepSpace$.
\latPK\ имеет шесть соседей, \latOCK --- восемь; обе реже и отпадают, раз требуется максимальная плотность при одной длине ребра
(\cref{fig:carrier-packing-sc,fig:carrier-packing-bcc,fig:carrier-packing-fcc}).

\begin{figure}[htbp]
  \centering
  \begin{tikzpicture}[carrier panel, iso view, scale=0.95]
  \def\r{0.38}
  \IsoSphereInk{0}{0}{0}{\r}
  \foreach \x/\y/\z in {1,0,0/-1,0,0/0,1,0/0,-1,0/0,0,1/0,0,-1} {
    \IsoSphereMuted{\x}{\y}{\z}{\r}
  }
  \node[anchor=north, font=\footnotesize] at (0,-1.15) {ПК, $|N|=6$};
\end{tikzpicture}

  \caption{\latPK\ (\latPKshort): $|N|=6$.}
  \label{fig:carrier-packing-sc}
\end{figure}

\begin{figure}[htbp]
  \centering
  \begin{tikzpicture}[carrier panel, iso view, scale=0.95]
  \def\r{0.36}
  \IsoSphereInk{0}{0}{0}{\r}
  \foreach \s in {-1,1} {
    \IsoSphereMuted{\s}{\s}{\s}{\r}
    \IsoSphereMuted{\s}{\s}{-\s}{\r}
    \IsoSphereMuted{\s}{-\s}{\s}{\r}
    \IsoSphereMuted{-\s}{\s}{\s}{\r}
  }
  \node[anchor=north, font=\footnotesize] at (0,-1.2) {BCC, $|N|=8$};
\end{tikzpicture}

  \caption{\latOCK\ (\latOCKshort): $|N|=8$.}
  \label{fig:carrier-packing-bcc}
\end{figure}

\begin{figure}[htbp]
  \centering
  \begin{tikzpicture}[carrier panel, iso view, scale=0.95]
  \def\r{0.34}
  \IsoSphereInk{0}{0}{0}{\r}
  \foreach \x/\y/\z in {
    1,1,0/1,-1,0/-1,1,0/-1,-1,0,
    1,0,1/1,0,-1/-1,0,1/-1,0,-1,
    0,1,1/0,1,-1/0,-1,1/0,-1,-1
  } {
    \IsoSphereMuted{\x}{\y}{\z}{\r}
  }
  \node[anchor=north, font=\footnotesize] at (0,-1.25) {FCC, $|N|=12$};
\end{tikzpicture}

  \caption{\latGCK\ (\latGCKshort): $|N|=12$.}
  \label{fig:carrier-packing-fcc}
\end{figure}

Плотность и $|N|$ у \latGCKshort\ и \latGPUshort\ совпадают; их различает симметрия.
\pfrom{требования изотропии вакуума}
у \latGCKshort\ три равные кубические оси, у \latGPUshort --- выделенная ось слоёв.
\phence носитель --- \latGCK.

\emph{Световой конус.}
Двенадцать соседей узла в начале координат, при ребре кубической ячейки $\caStepSpace\sqrt{2}$, суть все перестановки
\[
  \bigl(\pm\caStepSpace/\sqrt{2},\ \pm\caStepSpace/\sqrt{2},\ 0\bigr).
\]
Расстояние до начала
\[
  \sqrt{2\cdot(\caStepSpace/\sqrt{2})^2}=\caStepSpace.
\]
Вторая оболочка --- перестановки $(\pm\caStepSpace\sqrt{2},0,0)$, расстояние $\caStepSpace\sqrt{2}$.
\pusing{одного такта $\caSignalSpeed\htick=\caStepSpace$}
интервал до узла второй оболочки
\[
  \caSignalSpeed^2\htick^2-(\caStepSpace\sqrt{2})^2=\caStepSpace^2-2\caStepSpace^2=-\caStepSpace^2<0.
\]
\pfollows
узел второй оболочки за этот такт недостижим
(\cref{fig:carrier-light-cone-2d,fig:carrier-light-cone-3d}).
Остаётся первая оболочка: $|N(x)|=12$ (\cref{fig:carrier-fcc-shell}).

\begin{figure}[htbp]
  \centering
  \begin{tikzpicture}[carrier panel]
  \CarrierLightCone{1.15}{1.35}
  \fill[carrier muted, opacity=0.25] (-1.35,1.15) -- (0,0) -- (1.35,1.15) -- cycle;
  \node[carrier site] at (0,0) {};
  \node[carrier muted, font=\scriptsize] at (-0.35,0.55) {будущее};
  \draw[carrier object] (0,0) -- (1.0,1.15);
  \node[carrier site, minimum size=2.5pt] at (1.0,1.15) {};
  \draw[carrier muted, dashed] (0,0) -- (1.41,1.15);
  \node[carrier muted, font=\scriptsize] at (1.05,0.75) {$\caStepSpace\sqrt{2}$};
  \CarrierDimH{0}{1.0}{1.15}{$\caStepSpace$}
  \CarrierDimV{0}{1.15}{0}{$\caStepTime$}
\end{tikzpicture}

  \caption{Световой конус $\caSignalSpeed$ в срезе $(x,t)$: первая оболочка на конусе, вторая ($\caStepSpace\sqrt{2}$) --- вне.}
  \label{fig:carrier-light-cone-2d}
\end{figure}

\begin{figure}[htbp]
  \centering
  \begin{tikzpicture}[carrier panel, iso view, scale=0.95]
  \foreach \a in {0,45,...,315} {
    \draw[carrier muted] (0,0,0) -- ({1.1*cos(\a)},{1.1*sin(\a)},1.1);
  }
  \draw[carrier object] (0,0,0) -- (0,0,1.1);
  \fill[carrier fill, opacity=0.5] (0,0,0) -- (1.1,0,1.1) -- (0,1.1,1.1) -- cycle;
  \node[carrier site] at (0,0,0) {};
\end{tikzpicture}

  \caption{Световой конус в объёме $(x,y,t)$.}
  \label{fig:carrier-light-cone-3d}
\end{figure}

\begin{figure}[htbp]
  \centering
  \begin{tikzpicture}[carrier panel, iso view, scale=1.05]
  \def\r{0.36}
  \IsoSphereInk{0}{0}{0}{\r}
  \foreach \x/\y/\z in {
    1,1,0/1,-1,0/-1,1,0/-1,-1,0,
    1,0,1/1,0,-1/-1,0,1/-1,0,-1,
    0,1,1/0,1,-1/0,-1,1/0,-1,-1
  } {
    \IsoSphereMuted{\x}{\y}{\z}{\r}
  }
  \draw[carrier arrow, carrier muted] (0,0,0) -- (1,0,0);
  \node[carrier muted, font=\scriptsize, anchor=west] at (0.55,-0.15) {$\caStepSpace$};
\end{tikzpicture}

  \caption{Плотнейшая \latGCK: $|N|=12$ касающихся сфер.}
  \label{fig:carrier-fcc-shell}
\end{figure}

\emph{Связь.}
\pfrom{изотропии микрошага на двенадцати рёбрах}
\[
  \kappa_{\mathrm{link}}=\frac{1}{|N|}=\frac{1}{12}.
\]
\end{proof}

\subsection{Однотактовое тело: почему кубооктаэдр}
\label{sec:one-tick-hull}

Кубооктаэдр в тексте появляется не как «выбранный из каталога Архимеда'' многогранник,
а как \emph{имя} пространственного следа одного микрошага при уже принятом носителе.

\begin{definition}[Однотактовая оболочка на $\Mlayer$]
\label{def:one-tick-hull}
Пусть $N(x)=\{y:|y-x|=\caStepSpace\}$ --- первая координационная сфера (\cref{thm:fcc-choice}).
При условии~\ref{ax:A1} за один~$\htick$ из~$x$ достижимы только точки~$N(x)$.
Пространственное тело одного такта --- выпуклая оболочка $\mathcal{K}(x)=\mathrm{conv}\,N(x)$.
\end{definition}

На альтернативных решётках, отсечённых в доказательстве~\ref{thm:fcc-choice}, то же определение
даёт другие тела --- и это показывает, что многогранник \emph{следует} за решёткой, а не подменяет её:
\begin{center}
\small
\begin{tabular}{@{}lll@{}}
\toprule
решётка & $|N|$ & $\mathcal{K}(x)=\mathrm{conv}\,N(x)$ \\
\midrule
\latPKshort & $6$ & октаэдр \\
\latOCKshort & $8$ & куб \\
\latGCKshort\ / \latGPUshort & $12$ & кубооктаэдр / призма с шестиугольным основанием \\
\bottomrule
\end{tabular}
\end{center}
\latGPU\ с той же плотностью, что \latGCK, оставляет выделенную ось слоёв;
\pfrom{изотропии вакуума в \cref{thm:fcc-choice}}
носитель --- \latGCK, а не \latGPU.

\begin{proposition}[Кубооктаэдр как оболочка \latGCKshort]
\label{prop:cubocta-fcc}
Пусть $\carrierLattice$ --- \latGCK\ и $a=\caStepSpace$.
Тогда $N(0)$ --- все перестановки $(\pm a/\sqrt{2},\pm a/\sqrt{2},0)$,
и $\mathcal{K}(0)=\mathrm{conv}\,N(0)$ --- правильный кубооктаэдр с ребром $a/\sqrt{2}$,
шестью квадратными и восемью треугольными гранями (\cref{fig:carrier-cuboctahedron}).
\end{proposition}

\begin{proof}
Двенадцать точек лежат на сфере радиуса~$a$ и попарно симметричны относительно координатных плоскостей.
Любая грань $\mathrm{conv}\,N(0)$ обязана лежать в плоскости, несущей максимальную аффинную оболочку подмножества~$N(0)$.
Такие плоскости --- координатные квадраты (четыре соседа в одной плоскости $x=\mathrm{const}$)
и равносторонние треугольники на диагональных сечениях $x+y+z=\mathrm{const}$.
Других вершин у выпуклой оболочки нет: все точки $N(0)$ уже вершины.
Полученное тело совпадает с классическим кубооктаэдром (дуальный ромбическому додекаэдру --- ячейке Вороного на \latGCKshort).
\end{proof}

\begin{remark}[Что не рассматривается]
Икосаэдр, додекаэдр или усечённый октаэдр \emph{не} возникают как $\mathrm{conv}\,N(x)$ у периодической решётки
с одной длиной ребра~$\caStepSpace$ и первой координационной сферой в~$\mathbb{R}^3$:
у Браве-носителя соседи --- фиксированный дискретный набор, а не произвольные двенадцать точек на сфере.
Ячейка Вороного на \latGCKshort --- ромбический додекаэдр; это \emph{зона одного узла}, а не однотактовая каузальная оболочка (\cref{def:one-tick-hull}).
\end{remark}

\subsection{Геометрия однотактового тела на \latGCKshort}
\label{sec:fcc-kappa}

Далее метрика тела $\mathcal{K}(x)$ из \cref{prop:cubocta-fcc}.
Вершины --- те же точки $(\pm a/\sqrt{2},\pm a/\sqrt{2},0)$ и перестановки, где $a=\caStepSpace$.
Расстояние от центра до вершины уже сосчитано:
\[
  R_{\mathrm{out}}=a.
\]
Четыре точки $(a/\sqrt{2},\pm a/\sqrt{2},0)$ и $(a/\sqrt{2},0,\pm a/\sqrt{2})$ лежат в плоскости $x=a/\sqrt{2}$ и образуют квадрат со стороной~$a$.
Таких квадратных граней шесть.
Расстояние от центра до грани
\[
  R_{\mathrm{in}}^{\mathrm{hull}}=a/\sqrt{2}.
\]
Отношение двух длин одного и того же тела
\[
  \caKappaGeom=R_{\mathrm{in}}^{\mathrm{hull}}/R_{\mathrm{out}}=1/\sqrt{2}
\]
ещё не использует скорость света: это геометрия оболочки.
Треугольная грань через $(a/\sqrt{2},a/\sqrt{2},0)$, $(a/\sqrt{2},0,a/\sqrt{2})$ и $(0,a/\sqrt{2},a/\sqrt{2})$ лежит в плоскости $x+y+z=a\sqrt{2}$ и отстоит от центра на $a\sqrt{2/3}$.
Квадратные и треугольные грани лежат на разных расстояниях, поэтому одной вписанной сферы, касающейся всех граней, нет.
В мост $\speedC=\caKappaGeom \caSignalSpeed$ входит расстояние до квадратной грани
(\cref{fig:carrier-cuboctahedron,fig:carrier-cubocta-meridian}).

\begin{figure}[htbp]
  \centering
  \begin{tikzpicture}[carrier panel, iso view, scale=1.15]
  \foreach \x/\y/\z in {1,1,0/1,-1,0/-1,1,0/-1,-1,0/1,0,1/1,0,-1/-1,0,1/-1,0,-1/0,1,1/0,1,-1/0,-1,1/0,-1,-1} {
    \fill[carrier fill] (\x,\y,\z) circle (1.8pt);
  }
  \node[carrier site] at (0,0,0) {};
  \draw[carrier object] (1,1,0)--(1,0,1)--(0,1,1)--(-1,1,0)--(-1,0,1)--(0,-1,1)--(1,-1,0)--(1,0,-1)--(0,1,-1)--(-1,1,0);
  \draw[carrier muted, dashed] (1,1,0)--(0,1,1)--(-1,1,0);
  \fill[carrier hifill, opacity=0.5] (1,1,0)--(1,0,1)--(0,1,1)--cycle;
\end{tikzpicture}

  \caption{Кубооктаэдр --- выпуклая оболочка двенадцати соседей \latGCKshort.}
  \label{fig:carrier-cuboctahedron}
\end{figure}

\begin{figure}[htbp]
  \centering
  \begin{tikzpicture}[carrier panel]
  \draw[carrier muted] (-0.05,0) -- (1.3,0);
  \node[carrier site] at (0,0) {};
  \draw[carrier object] (0.72,0) -- (0.72,-0.38) -- (0.72,0.38);
  \draw[carrier arrow] (0,0) -- (0.72,0);
  \node[anchor=south, font=\scriptsize] at (0.36,-0.12) {$R_{\mathrm{in}}^{\mathrm{hull}}$};
  \draw[carrier arrow] (0,0) -- (0.55,0.48);
  \node[anchor=west, font=\scriptsize] at (0.6,0.5) {$R_{\mathrm{out}}=a$};
  \draw[carrier muted, dashed] (0.48,0.62) -- (0.72,0.62);
  \node[carrier muted, font=\scriptsize] at (0.6,0.75) {$a\sqrt{2/3}$};
\end{tikzpicture}

  \caption{Меридиан: квадратная и треугольная грани на разных расстояниях от центра.}
  \label{fig:carrier-cubocta-meridian}
\end{figure}

Объём ячейки Вороного считается отдельно.
Кубическая ячейка со стороной $A=a\sqrt{2}$ имеет объём $A^3=2\sqrt{2}\,a^3$ и содержит четыре узла \latGCKshort, поэтому на один узел
\[
  v_{\hv}=A^3/4=a^3/\sqrt{2}.
\]
Стена Вороного проходит посредине ребра к соседу, на расстоянии $a/2$ от узла.
Это не радиус оболочки:
\[
  R_{\mathrm{in}}^{\mathrm{Voronoi}}=a/2=\caKappaGeom\, R_{\mathrm{in}}^{\mathrm{hull}}.
\]
(\cref{fig:carrier-hull-voronoi,fig:carrier-voronoi-cell}).
Двенадцать соседей двойственны двенадцати ромбам этой ячейки.

\begin{figure}[htbp]
  \centering
  \begin{tikzpicture}[carrier panel, scale=1.45]
  \node[carrier site] at (0,0) {};
  \draw[carrier object] (0,0) circle (1);
  \draw[carrier muted] (0,0) circle ({1/sqrt(2)});
  \draw[carrier muted, dashed] (0,0) circle (0.5);
  \draw[carrier arrow] (0,0) -- (1,0);
  \node[anchor=south, inner sep=1pt] at (0.5,-0.06) {$R_{\mathrm{out}}=a$};
  \draw[carrier arrow, carrier muted] (0,0) -- ({0.5/sqrt(2)},{0.5/sqrt(2)});
  \node[carrier muted, font=\scriptsize, anchor=west] at (0.42,0.38) {$R_{\mathrm{in}}^{\mathrm{hull}}$};
  \draw[carrier arrow, carrier muted, dashed] (0,0) -- (0.5,0);
  \node[carrier muted, font=\scriptsize, anchor=north] at (0.5,0.08) {$R_{\mathrm{in}}^{\mathrm{Voronoi}}$};
  \node[anchor=north, font=\footnotesize] at (0,-1.25)
    {$\caKappaGeom=1/\sqrt{2}$};
\end{tikzpicture}

  \caption{Три радиуса однотактового тела на одном узле.}
  \label{fig:carrier-hull-voronoi}
\end{figure}

\begin{figure}[htbp]
  \centering
  \begin{tikzpicture}[carrier panel, iso view, scale=1.0]
  % schematic rhombic dodecahedron neighborhood
  \node[carrier site] at (0,0,0) {};
  \fill[carrier hifill, opacity=0.55]
    (1,0,0) -- (0,1,0) -- (-1,0,0) -- (0,-1,0) -- cycle;
  \draw[carrier object]
    (1,0,0) -- (0,1,0) -- (-1,0,0) -- (0,-1,0) -- cycle;
  \foreach \x/\y/\z in {1,1,0/1,-1,0/-1,1,0/-1,-1,0/1,0,1/1,0,-1/-1,0,1/-1,0,-1/0,1,1/0,1,-1/0,-1,1/0,-1,-1} {
    \draw[carrier muted, line width=0.4pt] (0,0,0) -- (\x*0.55,\y*0.55,\z*0.55);
    \fill[carrier fill, opacity=0.35] (\x*0.55,\y*0.55,\z*0.55) circle (1.5pt);
  }
  \draw[carrier object] (0.55,0.55,0) -- (0.55,0,0.55) -- (0,0.55,0.55) -- cycle;
\end{tikzpicture}

  \caption{Замощение $\mathbb{R}^3$ ячейками Вороного (ромбододекаэдры).}
  \label{fig:carrier-voronoi-cell}
\end{figure}

Объём самой оболочки находится вычитанием.
Многогранник с вершинами всех перестановок $(\pm 1,\pm 1,0)$ есть пересечение октаэдра $|x|+|y|+|z|\le 2$ и куба $|x|,|y|,|z|\le 1$.
Объём октаэдра $|x|+|y|+|z|\le 1$ равен $4/3$, поэтому при правой части~$2$ объём равен $(4/3)\cdot 8=32/3$.
Шесть шапок за гранями куба, по одной на полуось, имеют объём $2/3$ каждая: при $x$ от~$1$ до~$2$ сечение $|y|+|z|\le 2-x$ есть квадрат в повёрнутых осях площади $2(2-x)^2$, и
\[
  \int_1^2 2(2-x)^2\,\mathrm{d}x=\frac{2}{3}.
\]
Шесть шапок дают~$4$, остаётся $32/3-4=20/3$.
Ребро этого многогранника равно $\sqrt{2}$.
Масштаб к ребру~$a$ есть множитель $a/\sqrt{2}$, объём умножается на $(a/\sqrt{2})^3$:
\[
  V_{\mathrm{hull}}=\frac{20}{3}\cdot\frac{a^3}{2\sqrt{2}}=\frac{5\sqrt{2}}{3}\,a^3,
  \qquad
  \frac{V_{\mathrm{hull}}}{v_{\hv}}=\frac{5\sqrt{2}}{3}\cdot\sqrt{2}=\frac{10}{3}.
\]
Формула $(8\sqrt{2}/3)\,a^3$ с отношением $16/3$ --- объём другого тела.

Макроскопическая скорость связана с тактовой тем же отношением радиусов оболочки:
\begin{align}
  c &= \caKappaGeom \caSignalSpeed, \\
  \htick &= \caKappaGeom\, t_P, \\
  \caSignalSpeed &= \sqrt{2}\,c,
  \label{eq:hull-speed-bridge}
\end{align}
где $t_P=\caStepSpace/c$ --- учебниковая метка.
Непрерывное определение $t_P=\caStepSpace/c$ предполагало бы $\caSignalSpeed=c$; на этой оболочке тактовая и макроскопическая скорости различаются множителем~$\sqrt{2}$.

